% Ad Familiares 7.30 --- headnote
% ad-familiares-07-30, early January 44 BC, Rome

\textit{Cicero to M'. Curius at Patrae, written from Rome in the
first days of January 44 BC --- the reply to 7.29's request for a
commendation to Sulpicius's successor in Achaea. The Perseus
dateline gives \textit{in. m. Ian. a. 710 (44)}, ``the beginning of
January 44 BC,'' which matches the meta entry's month-precision
date.}

\textit{The substance is one of the most famous comic anecdotes in
the correspondence. On 31 December 45, the suffect consul Q. Fabius
Maximus dropped dead in the morning; Caesar, ever efficient, had a
fresh suffect (C. Caninius Rebilus) elected for the few remaining
hours of the year, so that Caninius's whole consulship was less than
half a day --- ``under the consulship of Caninius no man took his
luncheon.'' The sting is in the quotation of Ennius's lost
\textit{Iphigenia}: Cicero wants to flee somewhere he will not have
to hear of the doings of the house of Pelops --- that is, Caesar
and his circle. The closing sections turn back to Curius's affairs:
Cicero answers the property-law joke of 7.29 in kind (he is content
to own Curius by \textit{usus et fructus}), promises philosophy and
Atticus as his refuge from political ridicule, and encloses a
commendatory letter to the legate Acilius, who is in Curius's debt
to Cicero for two capital defences.}
